\providecommand{\Needspace}[1]{}
\Needspace{8\baselineskip}
\item \textbf{Clasificación de exoplanetas según el flujo recibido}.
\nopagebreak[4]

\begin{tcolorbox}[colback=blue!2!white,colframe=blue!55!black,title={Enunciado},breakable]
En astronomía, el flujo de radiación que recibe un exoplaneta depende de la luminosidad de su estrella y de la distancia orbital del planeta. En un modelo simplificado, este flujo puede estimarse mediante la expresión:
\[
S = \frac{L}{d^2},
\]
donde $S$ representa el flujo recibido, $L$ la luminosidad relativa de la estrella y $d$ la distancia orbital relativa del exoplaneta.

Considere los siguientes arreglos con información de distintos exoplanetas:
\[
L = [1.2,\ 0.8,\ 2.5,\ 1.6,\ 0.5,\ 3.0,\ 1.1],
\]
\[
d = [1.0,\ 0.7,\ 1.8,\ 1.2,\ 0.4,\ 2.5,\ 1.6].
\]

Cada elemento de ambos arreglos corresponde a un mismo exoplaneta. Es decir, la luminosidad \texttt{L[i]} y la distancia \texttt{d[i]} pertenecen al exoplaneta ubicado en el índice \texttt{i}.

Escriba un programa en Python que realice lo siguiente:

\begin{itemize}
    \item[(a)] Defina los arreglos de luminosidad $L$ y distancia $d$ usando \texttt{np.array}.

    \item[(b)] Cree una lista vacía donde se almacenarán los valores de flujo $S$ calculados para cada exoplaneta.

    \item[(c)] Recorra todos los índices de los arreglos usando un ciclo \texttt{for}. En cada iteración, acceda a \texttt{L[i]} y \texttt{d[i]} para calcular el flujo recibido por el exoplaneta correspondiente y agréguelo a la lista, por ejemplo con \texttt{lista\_vacia.append(calculo\_S)}:
    \[
    S_i = \frac{L_i}{d_i^2}.
    \]

    \item[(d)] Para cada exoplaneta, clasifique el flujo recibido usando condicionales según el siguiente criterio:
    \[
    S > 2 \quad \Longrightarrow \quad \text{Flujo alto},
    \]
    \[
    0.8 < S \leq 2 \quad \Longrightarrow \quad \text{Flujo intermedio},
    \]
    \[
    S \leq 0.8 \quad \Longrightarrow \quad \text{Flujo bajo}.
    \]

    \item[(e)] Muestre en pantalla, para cada exoplaneta, su índice, la luminosidad de su estrella, su distancia orbital, el flujo calculado y su clasificación.

    \item[(f)] Calcule e imprima el flujo promedio de todos los exoplanetas.

    \item[(g)] Realice un gráfico de dispersión usando \texttt{plt.scatter}, donde el eje horizontal represente la distancia $d$ y el eje vertical represente el flujo recibido $S$. Agregue título, nombres de ejes, grilla y una presentación limpia de la figura.
\end{itemize}

El objetivo del ejercicio es practicar el recorrido de arreglos mediante índices, usando expresiones como \texttt{L[i]} y \texttt{d[i]}, junto con condicionales para clasificar datos físicos y visualizar los resultados mediante un gráfico de dispersión.
\end{tcolorbox}

\begin{solution}
Una posible solución consiste en recorrer los arreglos con índices, calcular el flujo para cada exoplaneta, guardar los resultados en una lista y luego convertir esa lista a un arreglo de NumPy para calcular el promedio y graficar.

\begin{pythonjupyter}
import numpy as np
import matplotlib.pyplot as plt

# Arreglos de datos
L = np.array([1.2, 0.8, 2.5, 1.6, 0.5, 3.0, 1.1])
d = np.array([1.0, 0.7, 1.8, 1.2, 0.4, 2.5, 1.6])

# Lista vacía para guardar los flujos
flujos = []

# Recorrido usando índices
for i in range(len(L)):

    # Cálculo del flujo recibido
    S = L[i] / d[i]**2

    # Guardar el flujo calculado
    flujos.append(S)

    # Clasificación usando condicionales
    if S > 2:
        clasificacion = "Flujo alto"
    elif S > 0.8:
        clasificacion = "Flujo intermedio"
    else:
        clasificacion = "Flujo bajo"

    # Mostrar resultados de cada exoplaneta
    print("=" * 40)
    print("Exoplaneta índice:", i)
    print("Luminosidad:", L[i])
    print("Distancia:", d[i])
    print("Flujo:", S)
    print("Clasificación:", clasificacion)

# Convertir la lista de flujos a arreglo de NumPy
flujos = np.array(flujos)

# Flujo promedio
flujo_promedio = np.mean(flujos)

print("=" * 40)
print("Flujo promedio:", flujo_promedio)

# Gráfico de dispersión
plt.figure(figsize=(7, 5))
plt.scatter(d, flujos)

plt.title("Flujo recibido por exoplanetas")
plt.xlabel("Distancia orbital d")
plt.ylabel("Flujo recibido S")
plt.grid(True)

plt.show()
\end{pythonjupyter}
\end{solution}
